\begin{ExerciseBlock}
\ExerciseHeading{06}{Wellengleichung für das Magnetfeld}{wave_propagation_magnetic_wave_equation}
Im quellenfreien Vakuum gelten die Maxwell-Gleichungen
\[
\begin{aligned}
 \nabla\cdot\vec E&=0, & \nabla\cdot\vec B&=0,\\
 \nabla\times\vec E&=-\partial_t\vec B, &
 \nabla\times\vec B&=\mu_0\varepsilon_0\partial_t\vec E.
\end{aligned}
\]
Für das elektrische Feld folgt daraus $(\Delta-c^{-2}\partial_t^2)\vec E(\vec r,t)=0$. Leiten Sie die entsprechende Gleichung für die magnetische Flussdichte $\vec B$ her.
\begin{enumerate}
\item Überprüfen Sie für ein zweimal stetig differenzierbares Vektorfeld $\vec A(\vec r)$ die Identität
\[
 \nabla\times(\nabla\times\vec A)=\nabla(\nabla\cdot\vec A)-\Delta\vec A.
\]
\item Leiten Sie aus den Maxwell-Gleichungen im Vakuum die Gleichung $(\Delta-c^{-2}\partial_t^2)\vec B(\vec r,t)=0$ her.
\end{enumerate}
\end{ExerciseBlock}
\UseExerciseSolution{wave_propagation_magnetic_wave_equation}
